\scanpage{33}

\deleted{(iii) $\unsure{\longrightarrow|\Rightarrow}$ (i) (Cas $R$ réduit)}

\[
\text{(iii)}\underset{\text{(ii')}}{\Rightarrow}\left\{\begin{aligned}
&\exists\ \text{morphisme \deleted{\unsure{un}} fini surjectif net}\quad S'\to S,\ \text{tel que}\\
&\widehat A_{\widehat S'}\ \text{soit algébrique.}
\end{aligned}\right.
\]

Soit
\[
\begin{aligned}
&\deleted{\unsure{N_n=P_n|N_n\neq P_n}}\\
&N_n=\underline{NS}_{A_n/S_n}\qquad,\qquad
P_n=\text{\deleted{\unsure{partie}}schéma ouvert induit}\\
&\hspace{15em}\text{défini par les \uline{polarisations}}\\
&\hspace{13em}\text{\deleted{\unsure{ferm}}}
\end{aligned}
\]
Les $P_n$ se recollent et forment un schéma formel $\widehat P$ sur $\operatorname{Spf} R=\widehat S$, \deleted{\unsure{somme}} \inserted{provenant d'une famille \uuline{dénombrable}} de schémas finis nets $P^\alpha_{\deleted{\unsure{n}}}$ sur $S$
\[
\widehat P=\coprod_\alpha\widehat{P^\alpha}
\]

Soit
\[
Z_\alpha=\operatorname{Im}(P^\alpha\to S)
\]
L'hyp. (iii) signifie que pour tt $B$ comme dans l'énoncé, il existe un \inserted{$\alpha$ et un} $S$-morphisme \unsure{$\backslash$|!}
\[
\begin{aligned}
&\deleted{\operatorname{Spec} B'}\\
&T'\longrightarrow P^\alpha
\end{aligned}
\]
donc le pt générique $x$ de $T=\operatorname{Spec} B$ est dans $Z_\alpha$. Donc
\[
Z=\bigcup_\alpha Z_\alpha
\]
contient tous les $x\in S$ tels que $\dim\overline{\{x\}}=1$.

Par un lemme général, il s'ensuit qu'il y a une suite \uline{finie} de $Z_\alpha$ qui recouvrent $S$. Nous prenons pour $S'$ \deleted{\unsure{est}} la somme des $P^\alpha$ correspondants.

(ii') $\Rightarrow$ \deleted{(i) (Cas $R$ réduit)} \inserted{(i') :} $\widehat A_{\widehat S_{\unsure{1|3}}}$ est algébrique, si $S_1=S_{\text{réd}}$.

On peut supposer déjà $S$ réduit. Quitte à faire une extension étale finie $S_1\to S$ (et redescendre, \deleted{\unsure{$\exists$}} ce qui est possible par Raynaud), on peut supposer que
% archive sheet number: 32
